\section{Introduction}
\label{sec:introduction}

Instrumental-variables (IV) methods are a workhorse of empirical economics. Originally developed in the simultaneous-equations tradition of \cite{andersonEstimationParametersSingle1949} and refined for macroeconomic identification, the IV toolkit has, over the last three decades, migrated to applied microeconomics, where it underwrites a large empirical literature on labor markets, education, judges, examiners, and health \citep{angristTwoStageLeastSquares1995, angristInterpretationInstrumentalVariables2000, baiocchiTutorialBiostatisticsInstrumental2014}. Many of those applications appeal to multiple instruments to gain power: judge fixed effects, lottery-induced variation, shift-share decompositions, and Mendelian-randomization sets routinely involve dozens or hundreds of instruments. At the same time, applied work has moved away from the constant-effects model that motivated the classical IV asymptotics. Causal effects are presumed heterogeneous, and the two-stage least-squares (TSLS) estimand is interpreted as a positively-weighted average of local average treatment effects (LATEs), as in \cite{angristTwoStageLeastSquares1995}, \cite{kolesarEstimationInstrumentalVariables2013}, and \cite{mogstadCausalInterpretationTwoStage2021}.

This combination — many instruments, weak first stages, and heterogeneous treatment effects — is now the empirical default rather than the exception. Yet the workhorse weak-instrument-robust tests, against which applied researchers benchmark TSLS confidence intervals, were not designed for it. The \cite{andersonEstimationParametersSingle1949} AR test and its descendant, the \cite{moreiraConditionalLikelihoodRatio2003} CLR test, retain their nominal size under weak instruments only when treatment effects are constant. As soon as the data-generating process admits heterogeneous responses to the instruments — the very feature that justifies the LATE interpretation of TSLS — these tests over-reject, and the size distortion grows with instrument strength. The result is a gap in the inference toolkit precisely where applied work most often sits.

\paragraph*{Why the AR/CLR tests fail.}
The AR test inverts the moment restriction $\E[Z_i(Y_i - \beta_0 D_i)] = 0$, which is equivalent to the joint hypothesis that (a) there exists a scalar $\beta$ such that the reduced form satisfies $\delta_n = \beta\gamma_n$, and (b) that $\beta$ equals the hypothesized $\beta_0$. With more than one instrument, (a) is a substantive restriction: it is the constant-effects assumption in disguise. Treatment-effect heterogeneity violates (a), so the AR test rejects even when (b) is true. The decomposition $\mathrm{AR}(\beta_0) = \mathrm{LM}(\beta_0) + J(\beta_0)$ \citep{kleibergenGeneralizingWeakInstrument2007} makes this transparent: the $J$ component is a Sargan–Hansen test of overidentification \citep{sarganEstimationEconomicRelationships1958, hansenLargeSampleProperties1982}, which under strong exogeneity rejects whenever $\delta_n - \beta\gamma_n \neq 0$. The CLR test \citep{moreiraConditionalLikelihoodRatio2003} inherits the AR statistic through its construction and so inherits this failure mode. The LM test is robust to (a) but loses size under finite concentration parameters when heterogeneity is present.

\paragraph*{This paper.}
We propose a weak-instrument-robust test for the TSLS estimand $\beta = \gamma_n'\Sigma_{ZZ}\delta_n / \gamma_n'\Sigma_{ZZ}\gamma_n$ that retains its size under treatment-effect heterogeneity, weak instruments, and arbitrary numbers of instruments. The test is the asymptotic likelihood ratio for $H_0: \beta = \beta_0$, treating $H_0$ as a quadratic constraint on the reduced-form/first-stage coefficient pair $\tau_n \equiv (\delta_n', \gamma_n')'$. The constrained quasi-likelihood is a generalized trust-region subproblem in the sense of \cite{sternIndefiniteTrustRegion1995} and admits a closed-form solution at the optimum.

The resulting statistic, which we call TLR (TSLS likelihood ratio), has a tractable limiting distribution under strong exogeneity. The distribution depends on two scalar nuisance parameters: a magnitude-share $\varrho \in (0,1)$ governing the relative size of the positive and negative eigenvalues of the constraint matrix, and a non-centrality $\xi \geq 0$ summarizing instrument strength, treatment-effect heterogeneity, and endogeneity. We show that $\varrho$ has a consistent estimator and that $\xi$ admits a sufficient statistic with a chi-squared limit. With these in hand we construct two valid tests: a least-favorable-configuration (LFC) test that takes the worst-case quantile over $\xi$, and a Berger–Boos two-step test \citep{bergerValuesMaximizedConfidence1994} that first builds a confidence interval for $\xi$ and then takes the worst-case quantile over that interval. The Berger–Boos variant retains nontrivial power under both weak and strong instruments and, in our simulations, has size close to nominal and power close to the (size-distorted) TSLS Wald test that researchers currently use without correction.

The contribution is methodological: under standard heterogeneous-effects identification \citep{angristTwoStageLeastSquares1995, angristInterpretationInstrumentalVariables2000} and the weak-instrument asymptotic of \cite{staigerInstrumentalVariablesRegression1997}, our test is, to our knowledge, the first that delivers correct size for the TSLS estimand without requiring constant effects, a single instrument, or many-instrument asymptotics. The Section~\ref{sec:extensions} extension drops the strong-exogeneity restriction in favor of mean exogeneity and characterizes the limiting distribution in that setting; the construction continues to apply but with a larger nuisance space.

\paragraph*{Related literature.}
Our paper sits at the intersection of three literatures: (i) weak-instrument-robust inference under fixed instrument count, (ii) inference with many weak instruments, and (iii) identification under heterogeneous treatment effects.

The fixed-$d_Z$ weak-instrument literature begins with \cite{andersonEstimationParametersSingle1949} and was revived by \cite{staigerInstrumentalVariablesRegression1997}, \cite{boundProblemsInstrumentalVariables1995}, and the testing program of \cite{stockSurveyWeakInstruments2002} and \cite{stockTestingWeakInstruments2005}. The K-test of \cite{kleibergenGeneralizingWeakInstrument2007}, the CLR test of \cite{moreiraConditionalLikelihoodRatio2003}, the optimal invariant tests of \cite{andrewsOptimalTwoSidedInvariant2006}, and the practical implementations of \cite{finlayImplementingWeakInstrumentRobust2009} all maintain constant treatment effects. \cite{moreiraTestsCorrectSize2009} and \cite{andrewsConditionalLinearCombination2016} extend the conditional approach but, again, under constant effects. We share the conditional-inference perspective and the use of likelihood-ratio statistics, but we drop the constant-effects assumption and so live in a different limiting-distribution family — one indexed by $(\varrho, \xi)$ rather than by the single concentration parameter $\mu^2$.

The many-weak-instrument literature \citep{chaoConsistentEstimationLarge2005, chaoAsymptoticDistributionJive2012, mikushevaInferenceManyWeak2022, mikushevaWeakIdentificationMany2024, yapInferenceManyWeak2025} delivers tests that are valid under heterogeneous effects but only when $d_Z \to \infty$ along with $n$. The jackknife LM test of \cite{yapInferenceManyWeak2025} is the closest existing procedure to ours in spirit, though its asymptotics are taken in a different regime. Heteroskedasticity-robust inference with fixed $d_Z$ traces to \cite{kolesarEstimationInstrumentalVariables2013} and \cite{evdokimovInferenceInstrumentalVariables2018}; we maintain a heteroskedasticity-consistent variance estimator throughout.

The heterogeneous-effects identification literature begins with \cite{angristTwoStageLeastSquares1995}; the discrete-instrument LATE framework of % TODO[zotero]: Imbens \& Angrist (1994), "Identification and Estimation of Local Average Treatment Effects," \emph{Econometrica} 62(2): 467--475.
\cite{imbensAngristIdentificationEstimation1994} and the IV-as-LATE-aggregator argument of \cite{angristInterpretationInstrumentalVariables2000} and \cite{mogstadCausalInterpretationTwoStage2021} make explicit the conditions under which TSLS targets a meaningful causal estimand. % TODO[zotero]: Mogstad, Torgovitsky \& Walters (2021), "The Causal Interpretation of Two-Stage Least Squares with Multiple Instrumental Variables," \emph{American Economic Review} 111(11): 3663--3698. (Verify against the existing \cite{mogstadCausalInterpretationTwoStage2021} entry — they may be the same paper.)
% TODO[zotero]: Mogstad \& Torgovitsky (2025), \emph{Handbook of Econometrics, Volume X}, "Instrumental variables with heterogeneous treatment effects" (handbook chapter; consult for state-of-the-art positioning).
For the judge/examiner application that motivates much of our empirical interest, see  % TODO[zotero]: Chun \& Frandsen, "Examiner-design IV with heterogeneous treatment effects" (verify exact title and venue).
\cite{chunFrandsenExaminerDesign} and  % TODO[zotero]: Hull \& Goldsmith-Pinkham, "Examiner IV: Identification and inference" (verify exact title and venue).
\cite{hullGoldsmithPinkhamExaminer}, who document the prevalence of the multi-instrument, heterogeneous-effects regime in applied work.

\paragraph*{Roadmap.}
Section~\ref{sec:set-up-not} sets up the model and notation. Section~\ref{sec:ar-clr} reviews the AR, LM, and CLR tests and shows how each loses size under heterogeneity. Section~\ref{sec:results} introduces the TLR statistic, derives its limiting distribution, and constructs both the LFC and Berger–Boos tests. Section~\ref{sec:extensions} drops strong exogeneity. Section~\ref{sec:conclusion} concludes. Proofs are collected in Appendix~\ref{app:proofs}.
